\documentclass[11pt]{article}
\usepackage[margin=1in]{geometry}
\usepackage{amsmath,amssymb,amsthm,bm,xcolor,hyperref,booktabs,enumitem}
\newtheorem{theorem}{Theorem}
\newtheorem{proposition}{Proposition}
\newtheorem{corollary}{Corollary}
\newtheorem{definition}{Definition}
\newtheorem{assumption}{Assumption}
\newtheorem{conjecture}{Conjecture}
\newcommand{\PROVEN}[1]{{\color{teal}\small$\blacktriangleright$\,\textbf{[PROVEN/CHECKED]} #1}}
\newcommand{\SKETCH}[1]{{\color{orange}\small$\blacktriangleright$\,\textbf{[PROOF SKETCH]} #1}}
\newcommand{\CONJ}[1]{{\color{violet}\small$\blacktriangleright$\,\textbf{[TO PROVE]} #1}}
\newenvironment{intuition}{\par\smallskip\noindent\hangindent=1.5em\textbf{\itshape Intuition.}\itshape\ \ignorespaces}{\par\smallskip}
\newcommand{\R}{\mathbb{R}}\newcommand{\E}{\mathbb{E}}\newcommand{\N}{\mathcal{N}}
\renewcommand{\Pr}{\mathbb{P}}\newcommand{\I}{\mathbb{I}}\newcommand{\HH}{\mathbb{H}}

\title{\textbf{Structural De-aliasing of Latent States under Differential Misclassification}\\[3pt]
\large\normalfont(Paper B, redirected toward JASA T\&M --- theory work-through and roadmap)}
\author{Robert Richardson}
\date{\today}

\begin{document}
\maketitle

\begin{abstract}
A cheap ``oracle''---an LLM judge, an automated verifier, a weak annotator---measures a latent scientific state
$Z_v$ at each node $v$ of a known dependency structure, returning $Y_v$. We study the regime in which the
measurement channel \emph{aliases} two scientifically distinct states: $Y$ cannot separate them because their
node-level emission distributions coincide or nearly so. Our motivating pathology is \emph{bridge-blindness}: a
relevance judge grades a necessary \emph{bridge} passage exactly like an \emph{irrelevant} one, because it does
not directly answer the query. The question is not ``can we correct the judge?'' but the sharper statistical
one: \emph{when can the dependence structure distinguish latent states that the measurement channel itself
cannot?} We call this \textbf{structural de-aliasing}. We formalize a graph-indexed latent Markov model with a
node emission, define the aliasing depth $L^\star$, characterize exact aliasing by an emission-consistent
\emph{lumpable} partition (an impossibility boundary), separate \emph{functional} identifiability of a decision
target $R=h(Z)$ from full parameter identifiability, and give the information-theoretic limit
$\mathcal I_G=I(R_v;Y_{N(v)}\mid Y_v,X_v)$ that says exactly how much a cheap non-structural covariate leaves for
structure to add. We \emph{prove} the near-aliasing rate---a phase boundary $n(\delta_E^2+\delta_G^2)$ with
error exponent equal to the Chernoff information between the structural signatures, verified numerically---and lay
out the Bayesian weak-identification theory that would round out a T\&M paper. The empirical failure that prompted this
redirection---structure barely helps retrieval once cosine is present---is not a defeat: it is the regime
$\mathcal I_G\!\to\!0$, an instance the theory predicts exactly.
\end{abstract}

\paragraph{Status tags.} \PROVEN{} established or numerically verified here;\ \SKETCH{} argued, proof to be
completed;\ \CONJ{} the target results that make this T\&M, not yet proven.

\tableofcontents

\section{The problem, and why it is not the old one}
The previous conception of this paper asked whether corpus structure can \emph{correct} a bridge-blind LLM
relevance judge. Empirically it cannot do much in retrieval: a calibrated retriever prior recovers the judge's
missed bridges nearly as well as anything, and the graph adds little. That negative, however, exposes a deeper
and more general problem. The judge does not merely add noise; it makes two scientifically different states---a
\emph{necessary bridge} and an \emph{irrelevant} passage---produce the \emph{same} observation law. This is
\emph{differential misclassification} that is total on a pair of states: the measurement is \emph{aliased}.

\begin{intuition}
Ordinary measurement error blurs a signal; here the instrument \emph{collapses} two states onto one output
distribution. No amount of that instrument, or better calibration of it, separates them. The only remaining hope
is that the two states \emph{behave differently in the dependency structure}---a bridge sits next to direct
evidence, an irrelevant passage next to more irrelevant material---so that the \emph{neighbourhood} of
measurements distinguishes what the node measurement cannot. That is the question worth a theorem.
\end{intuition}

\paragraph{Positioning.} Generic identifiability of HMMs with linearly-independent emissions is classical
(Gassiat et al.); identifiability of HMMs with \emph{two aliased states} has been studied (Weiss \& Nadler,
2015); and a 2026 line proves nonidentifiability of LLM-judge debiasing and proposes anchor/paired designs. We
do \emph{not} re-enter those lanes. Our distinct contribution is the combination: \emph{(i)} graph/tree-indexed
latent structure, \emph{(ii)} \emph{heterophilic} compatibility (aliased states occur next to \emph{different}
states, not like themselves), \emph{(iii)} \emph{near}-aliasing rates rather than only generic identifiability,
\emph{(iv)} identifiability of a scientific \emph{functional} when the full model is not identified, and
\emph{(v)} the information-theoretic boundary that ties all of it to when a cheap covariate leaves anything for
structure to add.

\section{Model}
Let $G=(V,\mathcal E)$ be a known graph (we develop the tree/chain case first). Each node carries a latent state
$Z_v\in[K]=\{1,\dots,K\}$; on a tree the field is a Markov random field with edge transition matrix $T$
($T_{cc'}=\Pr(Z_u\!=\!c'\mid Z_v\!=\!c)$) and root/initial law $\pi$. A measurement channel (the oracle) emits
\begin{equation}
  Y_v \mid Z_v=c \;\sim\; O_c \quad(\text{row } c \text{ of the }K\times M\text{ emission matrix } O).
\end{equation}
Optionally each node has cheap non-structural covariates $X_v$ (in retrieval, the cosine prior). The scientific
target is a \emph{functional} $R_v=h(Z_v)$ (in retrieval, $h(c)=\I\{c\ \text{relevant}\}$). Two states $a,b$ are
\emph{node-aliased} if $O_a=O_b$ (exactly) or $d(O_a,O_b)=\delta_E$ small (near). Bridge-blindness is
$O_{\textsf{bridge}}=O_{\textsf{irrelevant}}$ with $O_{\textsf{direct}}$ distinct.

\section{Structural signatures and aliasing depth}
\begin{definition}[Structural observation signature]
$Q_L(c)=\mathcal L\big(Y_{B_L(v)}\mid Z_v=c\big)$, the law of the radius-$L$ neighbourhood of measurements given
the centre state. $Q_0(c)=O_c$.
\end{definition}
If the judge is perfectly bridge-blind, $Q_0(\textsf{bridge})=Q_0(\textsf{irrelevant})$. De-aliasing asks
whether $Q_L$ separates them for some $L>0$.

\begin{theorem}[De-aliasing and aliasing depth]\label{thm:dealias}
On a tree-indexed Markov chain, the law of the measurement at a node at path-distance $L$ from $v$, given
$Z_v=c$, is row $c$ of $T^{L}O$:
\[
   \Pr(Y_w=y\mid Z_v=c)=\big(T^{L}O\big)_{c,y}.
\]
Hence if $O_a=O_b$ but $(T^{L}O)_a\neq(T^{L}O)_b$ for some $L$, states $a,b$ are distinguishable from a
distance-$L$ measurement, and the \emph{aliasing depth} is
$L^\star(a,b)=\min\{L\ge0:(T^{L}O)_a\neq(T^{L}O)_b\}$.
\end{theorem}
\PROVEN{Proof is the Markov property: $\Pr(Y_w\mid Z_v{=}c)=\sum_{c'}\Pr(Y_w\mid Z_w{=}c')\,(T^{L})_{cc'}
=(T^{L}O)_{c}$. In the bridge-blind model with heterophilic $T$ (bridge$\to$direct), $O_{\textsf{irr}}=
O_{\textsf{bridge}}$ but $(TO)_{\textsf{irr}}=(0.64,0.20,0.16)\neq(TO)_{\textsf{bridge}}=(0.28,0.20,0.52)$, so
$L^\star=1$: one hop de-aliases. Verified: classifying the node-aliased pair from the node measurement alone is
at chance ($0.598$, prior-driven), from the neighbourhood $0.733$.}
\begin{intuition}
$L^\star$ is ``how many hops of context you need to tell the two states apart.'' Node-level aliasing is
$L^\star>0$; heterophily (aliased states have different neighbours) is what makes $L^\star$ finite and small.
\end{intuition}

\begin{theorem}[Exact aliasing $=$ lumpable partition; impossibility]\label{thm:lump}
Under a genericity condition on $O$, states $a,b$ are \emph{fully aliased} ($Q_L(a)=Q_L(b)$ for all $L$, hence
indistinguishable from any finite neighbourhood) iff they lie in the same block of the coarsest partition
$\mathcal P$ of $[K]$ that is (i) \emph{emission-consistent} ($O$ constant on blocks) and (ii) \emph{strongly
lumpable} for $T$ ($\sum_{c'\in B'}T_{cc'}$ is constant over $c\in B$, for all blocks $B,B'$). No consistent
estimator separates states in a common block.
\end{theorem}
\SKETCH{On a tree, conditioning on $Z_v=c$, the subtrees are independent and each is a Markov chain rooted at a
neighbour whose state is drawn from row $c$ of $T$; the whole neighbourhood law depends on $c$ only through that
row \emph{aggregated to the level at which the observation process is measurable}. Emission-consistency plus
strong lumpability (Kemeny--Snell) make the block process a Markov chain with block-measurable emissions, so
same-block states induce identical $Q_L$ for all $L$; conversely, if the pair is not in a common
emission-consistent lumpable block, some $(T^{L}O)$ row differs (Thm~\ref{thm:dealias}) and they de-alias. The
genericity condition rules out accidental cancellations in $O$. Making this fully rigorous---and extending
lumpability from chains to hidden Markov \emph{random fields}---is a core to-do.}

\section{Functional identifiability: decision recovery without parameter recovery}
The emission cannot see relevance ($O_{\textsf{bridge}}=O_{\textsf{irr}}$), so the full role model is not
node-identified. But the \emph{decision} may still be.
\begin{corollary}[Functional identifiability]\label{cor:func}
$R_v=h(Z_v)$ is identified from the neighbourhood iff every pair $a,b$ with $h(a)\neq h(b)$ is de-aliased
($L^\star(a,b)<\infty$). In the bridge-blind model the only $h$-splitting node-aliased pair is
$(\textsf{irrelevant},\textsf{bridge})$, de-aliased at radius $1$ iff $(TO)_{\textsf{irr}}\neq(TO)_{\textsf{bridge}}$
(bridges disproportionately adjacent to direct evidence). Hence \emph{relevance is recoverable even though the
emission is blind to it and the role parameters are not identified}.
\end{corollary}
\PROVEN{Immediate from Thm~\ref{thm:dealias}: identifying $h$ requires distinguishing only $h$-splitting pairs.}
This is the genuinely Bayesian/statistical distinction the paper should foreground: \emph{parameter recovery vs.\
decision recovery}. You need not identify the confusion matrix to identify what to do.

\section{The information limit: exactly when structure can help}
Let $X_v$ be cheap non-structural covariates. Define the \emph{incremental structural information}
\begin{equation}\label{eq:IG}
  \mathcal I_G \;=\; I\big(R_v;\,Y_{N(v)}\ \big|\ Y_v,\,X_v\big).
\end{equation}
\begin{theorem}[Information limit]\label{thm:info}
Under logarithmic loss, the reduction in optimal (Bayes) risk from adding the neighbourhood measurements is
exactly
\[
  \HH(R_v\mid Y_v,X_v)-\HH(R_v\mid Y_v,X_v,Y_{N(v)})=\mathcal I_G .
\]
In particular $\mathcal I_G=0\Rightarrow$ \emph{no} graph-based rule---however elaborate the MRF---improves optimal
log-loss prediction of $R_v$; and by Fano/Pinsker the improvement in $0$--$1$ Bayes risk is bounded by an
increasing function of $\mathcal I_G$.
\end{theorem}
\PROVEN{The identity is the definition of conditional mutual information. The $0$--$1$ bound follows from
Fano/Pinsker. \textbf{Verified regime-switch:} with $R=\I\{Z\ge1\}$ and a Gaussian covariate $X_v\sim\N(\mu_{Z_v},
\sigma)$ that separates relevant from irrelevant, the neighbour-gain (an $\mathcal I_G$ proxy in nats) is
$0.09$ (no $X$), $0.11$ (weak, $\sigma{=}1$), $0.10$ ($\sigma{=}0.5$), $0.007$ (strong, $\sigma{=}0.2$),
$0.000$ (near-perfect, $\sigma{=}0.1$): as the covariate becomes sufficient for $R$, $\mathcal I_G\to0$.}
\begin{intuition}
This single quantity explains all three things we saw. In the clean simulation there was no cheap covariate, so
$\mathcal I_G$ was large and structure rescued relevance. In real retrieval the cosine prior is a \emph{strong}
covariate---it already separates buried-bridge from distractor---so $\mathcal I_G$ is tiny and the graph adds
almost nothing. And when the covariate is sufficient, $\mathcal I_G=0$ and \emph{no} method can help. The
retrieval ``failure'' is the theorem's $\mathcal I_G\!\to\!0$ corner, not a bug.
\end{intuition}

\section{The near-aliasing rate (proved) and Bayesian weak identification}
The results above are qualitative (identified / not). T\&M wants \emph{rates and limits} in the near-aliased
regime, $\delta_E^2=H^2(O_a,O_b)\to0$ and $\delta_G^2=H^2\big((T^{L}O)_a,(T^{L}O)_b\big)\to0$
(squared Hellinger distances). The conditional-independence structure of the model makes this a clean product-of-channels
computation, and the conjectured phase boundary becomes a theorem.

\begin{theorem}[Near-aliasing rate and phase boundary]\label{thm:rate}
Test $H_a\!:Z_v=a$ against $H_b\!:Z_v=b$ (with $h(a)\neq h(b)$) from the node measurement $Y_v$ and $d$
neighbour measurements at radius $L$, under the tree-indexed model. Let the node and per-neighbour Chernoff
informations be $C_E=\max_{s}-\log\sum_y O_{a,y}^{s}O_{b,y}^{1-s}$ and
$C_G=\max_{s}-\log\sum_y (T^{L}O)_{a,y}^{s}(T^{L}O)_{b,y}^{1-s}$, and the Bhattacharyya coefficients
$\rho_E=\sum_y\sqrt{O_{a,y}O_{b,y}}=1-\tfrac12\delta_E^2$, $\rho_G=1-\tfrac12\delta_G^2$. Then for the
equal-prior Bayes error $P_e$:
\begin{enumerate}[label=(\alph*),leftmargin=*,itemsep=0pt]
\item \textbf{(product factorization)} conditional independence gives the joint law
$P_c=O_c\otimes\bigotimes_{j=1}^{d}(T^{L}O)_c$ under $H_c$, so the affinity multiplies:
$\rho(P_a,P_b)=\rho_E\,\rho_G^{\,d}$;
\item \textbf{(bounds)} $\tfrac14\big(\rho_E\rho_G^{d}\big)^{2}\le P_e\le \tfrac12\,\rho_E\rho_G^{d}$, hence
$P_e=\exp\!\big(-\Theta(\delta_E^2+d\,\delta_G^2)\big)$;
\item \textbf{(sharp exponent)} $P_e\;\asymp\;d^{-1/2}\exp\!\big(-(C_E+d\,C_G)\big)$ (Bahadur--Rao);
\item \textbf{(phase boundary)} $P_e\to0$ iff $\delta_E^2+d\,\delta_G^2\to\infty$; in the near-aliased regime with
$\delta_E$ bounded, iff $d\,\delta_G^2\to\infty$. Writing $n$ for the effective count of informative neighbours,
the boundary is $n(\delta_E^2+\delta_G^2)\to\infty$.
\end{enumerate}
\end{theorem}
\PROVEN{\emph{Proof.} (a) By the Markov property on the tree, given $Z_v=c$ the node emission and the $d$
neighbour subtrees are mutually independent, each neighbour measurement drawn from row $c$ of $T^{L}O$; the joint
law is the stated product. Hellinger affinity is multiplicative over product measures, giving
$\rho(P_a,P_b)=\rho_E\rho_G^{d}$. (b) Equal-prior Bayes error obeys the Bhattacharyya upper bound
$P_e\le\tfrac12\rho(P_a,P_b)$ and the Le~Cam lower bound $P_e\ge\tfrac12\big(1-\mathrm{TV}\big)\ge
\tfrac12\big(1-\sqrt{1-\rho^2}\big)\ge\tfrac14\rho^{2}$; substitute $\rho=\rho_E\rho_G^{d}$ and use
$-\log\rho_\bullet=\tfrac12\delta_\bullet^2(1+o(1))$. (c) The log-likelihood ratio is a sum of $d$ i.i.d.\ neighbour
terms plus the node term; Chernoff's theorem gives error exponent equal to the summed Chernoff informations
$C_E+dC_G$, and the local CLT supplies the $d^{-1/2}$ prefactor (Bahadur--Rao). (d) Immediate from (b)--(c).
\hfill$\square$\\[2pt]
\textbf{Verified} (\texttt{paperB\_rate\_sim.py}, exact node alias $\delta_E{=}0$, so all information is
structural): with $(T\!O)_a{=}(.64,.20,.16)$, $(T\!O)_b{=}(.28,.20,.52)$ the Chernoff information is $C_G{=}0.0925$
($s^\star{=}0.49$, so Bhattacharyya $\beta_G{=}0.0924$ nearly attains it). A Bahadur--Rao fit
$-\log P_e = a\,d + p\log d + \text{const}$ over $d\le 96$ recovers exponent $a=0.0934$ (matches $C_G=0.0925$) and
polynomial order $p=0.34$ (matches the predicted $\tfrac12$): the error decays as $d^{-1/2}e^{-dC_G}$ with the
\emph{structural} channel alone, though the measurement is fully aliased.}
\begin{intuition}
The rate says exactly how ``buy-able'' de-aliasing is: each informative neighbour multiplies the odds of the
right call by a fixed factor $e^{C_G}$, so error falls geometrically in the number of neighbours---but the
per-neighbour exponent $C_G\approx\tfrac14\delta_G^2$ is tiny when the states are near-aliased, so you need
$\sim 1/\delta_G^2$ neighbours to separate them. That $1/\delta_G^2$ is the price of the alias.
\end{intuition}
\paragraph{Remaining rigor.} (a)--(d) are proved for the immediate-neighbour (degree-$d$ / star) case, which is
also the retrieval geometry. Two extensions complete the T\&M version: (i) \emph{general trees}---the affinity
still factorizes over the $d$ immediate subtrees and each subtree affinity is $\le$ its radius-$1$ value, so the
bound holds a fortiori; the exact deep-subtree Chernoff constant is the remaining computation; (ii) the
\emph{field} minimax (estimating $R$ over all nodes) follows from the pairwise rate by a standard Assouad/Fano
reduction. Neither is expected to change the boundary $n(\delta_E^2+\delta_G^2)$.
\begin{conjecture}[Bayesian weak identification]\label{conj:bayes}
As $(\delta_E,\delta_G)\to0$: (a) the posterior for $R$ contracts iff $n(\delta_E^2+\delta_G^2)\to\infty$; (b)
the emission-bias nuisance may fail to contract while $R$ does (functional vs.\ parameter); (c) two admissible
structural priors $\Pi_1,\Pi_2$ give $d_{\mathrm{TV}}\{\Pi_1(R\mid Y,G),\Pi_2(R\mid Y,G)\}\to0$ in the identified
regime and bounded-away-from-zero in the structurally-unidentified regime---a computable diagnostic separating
``a correction learned from data'' from ``a correction imposed by the structural prior.''
\end{conjecture}
\CONJ{Posterior-contraction machinery (Ghosal--van der Vaart style) adapted to the tree-signature likelihood;
the prior-sensitivity claim is a cut/robustness statement. This is the most Bayesian-Analysis-flavoured piece and
is the natural second paper if the T\&M rate theory is the first.}

\section{The three empirical regimes as instances of one theory}
\begin{center}\small
\begin{tabular}{lccl}
\toprule
regime & cheap covariate $X$ & $\mathcal I_G$ & outcome \\ \midrule
favorable simulation (Paper B v1) & absent/weak & large & graph rescues relevance \\
real retrieval (Paper B v1) & cosine (strong) & $\approx 0$ & graph adds $\approx\!+0.03$; prior does the work \\
impossibility & sufficient for $R$ & $0$ & no method can help \\ \bottomrule
\end{tabular}
\end{center}
The apparent contradiction between the favorable sim and the retrieval null dissolves: they are two points on the
$\mathcal I_G$ axis of Theorem~\ref{thm:info}.

\section{Roadmap to a JASA T\&M paper}
\begin{enumerate}[leftmargin=*,itemsep=1pt]
  \item \textbf{Finish Thms \ref{thm:dealias}--\ref{thm:info}} rigorously (\ref{thm:dealias},\ref{cor:func},
  \ref{thm:info} essentially done; \ref{thm:lump} needs the lumpability$\to$HMRF extension and the genericity
  condition).
  \item \textbf{Near-aliasing rate (Thm~\ref{thm:rate})---proved for the immediate-neighbour case} and
  numerically confirmed (exponent $=$ Chernoff information, $d^{-1/2}$ prefactor). Remaining: the deep-subtree
  Chernoff constant on general trees, and the Assouad/Fano reduction for the field minimax. This is the T\&M
  centerpiece---the phase boundary $n(\delta_E^2+\delta_G^2)$ is now a theorem, not a conjecture.
  \item \textbf{Measurement-design corollary.} If a second channel $O^{(2)}$ (a different prompt/model/verifier)
  partially breaks the alias, derive the \emph{structure$\leftrightarrow$replication trade-off}: how much
  targeted re-measurement restores identification as a function of $\delta_G$. (Positions us away from the 2026
  anchor/paired-design debiasing work: our knob is \emph{structure}, replication is the corollary.)
  \item \textbf{Screen a second, no-cheap-prior application by $\mathcal I_G$ before modelling.} Compute
  $\widehat{\mathcal I_G}=\hat I(R;\text{neighbour outputs}\mid \text{node output},\text{all cheap features})$ on
  candidate domains---verification of reasoning steps in a dependency DAG, program/SQL verification on an
  AST/edit graph, claim/evidence checking in a RAG chain---and \emph{only} build the model where $\mathcal I_G$
  is materially positive after every cheap predictor. (Retrieval taught us this screen is mandatory.)
  \item \textbf{Bayesian companion} (Conj.~\ref{conj:bayes}) for a Bayesian Analysis version: posterior
  contraction, prior sensitivity, and the ``learned vs.\ imposed correction'' diagnostic.
\end{enumerate}

\paragraph{Ceiling assessment.} A chain-Kruskal identification theorem plus the LLM experiment is below T\&M and
too narrow for BA. \emph{Structural de-aliasing $+$ the near-aliasing rate $+$ functional Bayes} is a legitimate
JASA T\&M shot; \emph{de-aliasing $+$ posterior-contraction / prior-sensitivity} is a legitimate Bayesian
Analysis shot. The failure we found did not kill B---it told us the original theorem was aimed at the easy part
of the parameter space, and handed us the boundary that is the interesting part.

\end{document}
